% Teil II: der durchgerechnete Weg.

\subsection*{Gegeben und gesucht}
Beweisen: Die Summe der Innenwinkel eines $n$-Ecks ($n\ge3$, Eckpunkte
auf einer gemeinsamen Kreislinie) ist $(n-2)\cdot180^\circ$.

\subsection*{Ansatz}
Vollständige Induktion über die Eckenzahl $n$, im Schritt per Diagonale
zerlegen.

\subsection*{Induktionsanfang}
$n=3$: Dreieckssumme ist $180^\circ=(3-2)\cdot180^\circ$. \checkmark

\subsection*{Induktionsschritt}
Sei die Aussage für ein $n$-Eck bewiesen. Ein $(n+1)$-Eck (konvex) wird
durch eine Diagonale in ein Dreieck und ein $n$-Eck zerlegt. Die
Innenwinkelsumme des $(n+1)$-Ecks ist die des $n$-Ecks plus die des
Dreiecks:
\begin{equation}
  S(n+1)=S(n)+180^\circ=(n-2)180^\circ+180^\circ=(n-1)180^\circ .
\end{equation}
Wegen $n-1=(n+1)-2$ ist das die Behauptung für $n+1$.

\begin{ergebnis}
\begin{equation*}
  S(n)=(n-2)\cdot180^\circ\quad\text{für alle } n\ge3 .
\end{equation*}
\end{ergebnis}

\begin{kontrolle}
$n=4$: Viereck hat $\text{360}^\circ=(4-2)180^\circ$; $n=5$: Fünfeck
$\text{540}^\circ=(5-2)180^\circ$. Vertraut und plausibel.
\end{kontrolle}

\begin{merksatz}
Ein konvexes Polygon zerlegt man in Dreieck plus kleineres Polygon; jede
Zerlegung addiert $180^\circ$, was auf die Formel $(n-2)\cdot180^\circ$
führt.
\end{merksatz}
